\setbeamertemplate{note page}[plain]

\lstset{ % General setup for the package
    language=Java, 
    basicstyle=\scriptsize\sffamily,
    numbers=left,
    numberstyle=\tiny, 
    frame=tb,
    tabsize=2,
    columns=fixed,
    showstringspaces=false,
    showtabs=false,
    keepspaces,
    commentstyle=\color{red},
    keywordstyle=\color{blue}
}



\newcommand{\VORLESUNG}{1}
\newcommand{\KURSEINHEITNR}{36}
\newcommand{\NAMEVORLESUNG}{Datenbanken Repetitorium}

\if \VORLESUNG 1
    \setbeameroption{}
\else 
    \setbeameroption{show notes on second screen}
\fi 

\date{Datenbanken Repetitorium}
\newcommand{\BILDERPFAD}{\KURSEINHEIT/bilder}

\ifnum\KURSEINHEITNR=0
    \newcommand{\KURSEINHEIT}{00-Example-Chapter}
    \newcommand{\presentationTitle}{Latex Beamer Template}
\fi

\ifnum\KURSEINHEITNR=1
    \newcommand{\KURSEINHEIT}{1-Einfuehrung}
    \newcommand{\presentationTitle}{Einführung}
\fi
	
\ifnum\KURSEINHEITNR=2
    \newcommand{\KURSEINHEIT}{2-UML}
    \newcommand{\presentationTitle}{Einführung in UML}
\fi

\ifthenelse {\KURSEINHEITNR = 3} {
    \newcommand{\KURSEINHEIT}{3-Singleton}
    \newcommand{\presentationTitle}{Singleton}
}{}

\ifthenelse {\KURSEINHEITNR = 4} {
    \newcommand{\KURSEINHEIT}{4-Multiton}
    \newcommand{\presentationTitle}{Multiton}
}{}

\ifthenelse {\KURSEINHEITNR = 5} {
    \newcommand{\KURSEINHEIT}{5-Decorator}
    \newcommand{\presentationTitle}{Dekorator}
}{}

\ifthenelse {\KURSEINHEITNR = 6} {
    \newcommand{\KURSEINHEIT}{6-Strategie}
    \newcommand{\presentationTitle}{Strategie}
}{}

\ifthenelse {\KURSEINHEITNR = 7} {
    \newcommand{\KURSEINHEIT}{7-Schablonenmethode}
    \newcommand{\presentationTitle}{Schablonenmethode}
}{}

\ifthenelse {\KURSEINHEITNR = 8} {
    \newcommand{\KURSEINHEIT}{8-kompositum}
    \newcommand{\presentationTitle}{Kompositum}
}{}

\ifthenelse {\KURSEINHEITNR = 9} {
    \newcommand{\KURSEINHEIT}{9-Iterator}
    \newcommand{\presentationTitle}{Iterator}
}{}

\ifthenelse {\KURSEINHEITNR = 10} {
    \newcommand{\KURSEINHEIT}{10-Bruecke}
    \newcommand{\presentationTitle}{Brücke}
}{}

\ifthenelse {\KURSEINHEITNR = 11} {
    \newcommand{\KURSEINHEIT}{11-Kommando}
    \newcommand{\presentationTitle}{Kommando}
}{}

\ifthenelse {\KURSEINHEITNR = 12} {
    \newcommand{\KURSEINHEIT}{12-Proxy}
    \newcommand{\presentationTitle}{Proxy}
}{}

\ifthenelse {\KURSEINHEITNR = 13} {
    \newcommand{\KURSEINHEIT}{13-Fabrikmethode}
    \newcommand{\presentationTitle}{Fabrikmethode}
}{}

\ifthenelse {\KURSEINHEITNR = 14} {
    \newcommand{\KURSEINHEIT}{14-AbstrakteFabrik}
    \newcommand{\presentationTitle}{Abstrakte Fabrik}
}{}

\ifthenelse {\KURSEINHEITNR = 15} {
    \newcommand{\KURSEINHEIT}{15-Prototyp}
    \newcommand{\presentationTitle}{Prototyp}
}{}

\ifthenelse {\KURSEINHEITNR = 16} {
    \newcommand{\KURSEINHEIT}{16-Mediator}
    \newcommand{\presentationTitle}{Mediator}
}{}

\ifthenelse {\KURSEINHEITNR = 17} {
    \newcommand{\KURSEINHEIT}{17-Fliegengewicht}
    \newcommand{\presentationTitle}{Fliegengewicht}
}{}

\ifthenelse {\KURSEINHEITNR = 18} {
    \newcommand{\KURSEINHEIT}{18-Beobachter}
    \newcommand{\presentationTitle}{Beobachter}
}{}

\ifthenelse {\KURSEINHEITNR = 19} {
    \newcommand{\KURSEINHEIT}{19-Erbauer}
    \newcommand{\presentationTitle}{Erbauer}
}{}

\ifthenelse {\KURSEINHEITNR = 20} {
    \newcommand{\KURSEINHEIT}{20-Besucher}
    \newcommand{\presentationTitle}{Besucher}
}{}

\ifthenelse {\KURSEINHEITNR = 21} {
    \newcommand{\KURSEINHEIT}{21-Fassade}
    \newcommand{\presentationTitle}{Fassade}
}{}

\ifthenelse {\KURSEINHEITNR = 22} {
    \newcommand{\KURSEINHEIT}{22-Adapter}
    \newcommand{\presentationTitle}{Adapter}
}{}

\ifthenelse {\KURSEINHEITNR = 23} {
    \newcommand{\KURSEINHEIT}{23-Zustand}
    \newcommand{\presentationTitle}{Zustand}
}{}

\ifthenelse {\KURSEINHEITNR = 24} {
    \newcommand{\KURSEINHEIT}{24-ChainOfResponsibility}
    \newcommand{\presentationTitle}{Verantwortlichkeitskette}
}{}

\ifthenelse {\KURSEINHEITNR = 25} {
    \newcommand{\KURSEINHEIT}{25-Memento}
    \newcommand{\presentationTitle}{Memento}
}{}

\ifthenelse {\KURSEINHEITNR = 26} {
    \newcommand{\KURSEINHEIT}{26-Interpreter}
    \newcommand{\presentationTitle}{Interpreter}
}{}

\ifthenelse {\KURSEINHEITNR = 27} {
    \newcommand{\KURSEINHEIT}{27-DependencyInjection}
    \newcommand{\presentationTitle}{Dependency Injection}
}{}

\ifthenelse {\KURSEINHEITNR = 28} {
    \newcommand{\KURSEINHEIT}{28-FluentAPI}
    \newcommand{\presentationTitle}{FluentAPI}
}{}

\ifthenelse {\KURSEINHEITNR = 29} {
    \newcommand{\KURSEINHEIT}{29-Immutable}
    \newcommand{\presentationTitle}{Unveränderliche Objekte}
}{}

\ifthenelse {\KURSEINHEITNR = 30} {
    \newcommand{\KURSEINHEIT}{30-Specification}
    \newcommand{\presentationTitle}{Specification Pattern}
}{}

\ifthenelse {\KURSEINHEITNR = 31} {
    \newcommand{\KURSEINHEIT}{31-NullObjectPattern}
    \newcommand{\presentationTitle}{Null Object Pattern}
}{}

\ifthenelse {\KURSEINHEITNR = 32} {
    \newcommand{\KURSEINHEIT}{32-Repository}
    \newcommand{\presentationTitle}{Repository}
}{}

\ifthenelse {\KURSEINHEITNR = 33} {
    \newcommand{\KURSEINHEIT}{33-ValueObjects}
    \newcommand{\presentationTitle}{Value Objects}
}{}

\ifthenelse {\KURSEINHEITNR = 34} {
    \newcommand{\KURSEINHEIT}{34-ModelViewController}
    \newcommand{\presentationTitle}{ModelViewController}
}{}

\ifthenelse {\KURSEINHEITNR = 35} {
    \newcommand{\KURSEINHEIT}{35-Einführung-Datenbanken}
    \newcommand{\presentationTitle}{Einführung}
}{}

\ifthenelse {\KURSEINHEITNR = 36} {
    \newcommand{\KURSEINHEIT}{36-Datenbankentwurf}
    \newcommand{\presentationTitle}{Datenbankentwurf}
}{}